\subsection{Energie théorique maximale}


\setcounter{equation}{0}
\setcounter{Q}{0}

A l'exercice précédent, nous avons vu qu'il était possible de calculer les performances théoriques maximales d'un matériau d'électrode.

En particulier, ces performances sont mesurées à l'aide d'un accumuluteur présentant un pôle $\ominus$ (souvent appelé anode dans la littérature) en lithium métallique : 

\begin{align*}
\oplus \quad Li_x M(s) | Li^+ || Li^+ | Li(s) \quad \ominus
\end{align*}

où $M$ désigne un matériau d'életrode donné.

Afin de comparer les performances des électrodes du pôle $\oplus$ (souvent appelé cathode dans la littérature), nous allons proposer les potentiels suivants : 

\begin{center}
\begin{tabular}{ccccccc}
\toprule
Matériau & \ce{Li(s)} & \ce{LiC6(s)} & \ce{Li15Si4(s)} & \ce{LiCoO2(s)} & \ce{LiFePO4(s)} & \ce{LiNi_{1/3}Co_{1/3}Mn_{1/3}O2(s)} \\
\midrule
$E_{\oplus}^{eq}$ [$\si{\volt\per{Li^+/Li}}$] & 0 & 0.3 & 0.6 & 4.1 & 3.3 & 3.9  \\
\bottomrule
\end{tabular}
\captionof{figure}{Potentiel électrochimique d'équilibre à \SI{298}{\kelvin} de différents matériaux du pôle $\oplus$}
\end{center}



\Q Qu'appelle-t-on une pseudo-référence ? Qui joue ce rôle ici ? Que signifie $\si{\volt\per{Li^+/Li}}$ ?

\solution{
Les électrodes de pseudo-référence sont des électrodes fabriquée à partir de matériaux dont le potentiel électrochimique à l'aquilibre est stable ou dont la variation est prévisible.\\
L'électrode au lithium métallique est considérée ici comme une pseudo-référence dont le potentiel est supposé constant et valoir approximativement : 
\begin{align*}
E_\oplus^{eq} \approx - 3.04 ~ \si{\volt\per{ESH}}
\end{align*}
L'unité $\si{\volt\per{Li^+/Li}}$ signifie que l'on considère le potentiel de l'électrode $\oplus$ par rapport au couple \ce{Li^+}/\ce{Li} à l'électrode $\ominus$ : 
\begin{align*}
E_\oplus^{eq} [\si{\volt\per{Li^+/Li}}] &= E_\oplus^{eq} [\si{\volt\per{ESH}}] - E_\ominus^{eq} [\si{\volt\per{ESH}}] \\
&=E_\oplus^{eq} [\si{\volt\per{ESH}}] - E^{eq} (\ce{Li^+} / \ce{Li})
\end{align*}
}

\Q En déduire la f.e.m. dans ces cellules Li-metal. 
\solution{Par conséquent, la force électromotrice dans ces cellules est approximativement égale au potentiel du pôle $\oplus$ :
\begin{center}
\begin{tabular}{ccccccc}
\toprule
Matériau & \ce{Li(s)} & \ce{LiC6(s)} & \ce{Li15Si4(s)} & \ce{LiCoO2(s)} & \ce{LiFePO4(s)} & \ce{LiNi_{1/3}Co_{1/3}Mn_{1/3}O2(s)} \\
\midrule
$f.e.m.$ [$\si{\volt}$] & 0 & 0.3 & 0.6 & 4.1 & 3.3 & 3.9  \\
\bottomrule
\end{tabular}
\end{center}
}

\Q On observe deux catégories de matériaux à partir des valeurs de f.e.m. Quelles sont-elles ?  Donner un sens énergétique à la f.e.m.
\solution{On peut classer les matériaux en fonction de leur proximité en potentiel avec l'électrode au lithium. Les électrodes $\ominus$ (anode en décharge) doivent avoir un potentiel le plus proche possible du lithium, tandis que les électrodes $\oplus$ (cathode en décharge) doivent avoir le potentiel le plus éloigné possible du lithium métallique : 
\begin{center}
\begin{tabular}{ccccccc}
\toprule
Matériau & \ce{Li(s)} & \ce{LiC6(s)} & \ce{Li15Si4(s)} & \ce{LiCoO2(s)} & \ce{LiFePO4(s)} & \ce{LiNi_{1/3}Co_{1/3}Mn_{1/3}O2(s)} \\
\midrule
$f.e.m.$ [$\si{\volt}$] & 0 & 0.3 & 0.6 & 4.1 & 3.3 & 3.9  \\
Rôle & $\ominus$ & $\ominus$ & $\ominus$ & $\oplus$ & $\oplus$ & $\oplus$ \\
\bottomrule
\end{tabular}
\end{center}
} 

\Q Après avoir exprimé le travail électrique, calculer l'énergie théorique maximale stockable dans ces cathodes.
\solution{
En convention récepteur (resp. générateur), la puissance électrique reçue (resp. cédée) vaut : 
\begin{align*}
P = U \times I
\end{align*}
Le travail électrique reçu vaut : 
\begin{align*}
W' = \int P(t) dt = \int U(Q) dQ
\end{align*}
Au maximum, en décharge, la tension vaut la force électromotrice. L'énergie théorique maximale s'écrit donc :
\begin{align*}
ETM [\si{\milli\watt\hour\per\gram}] = e [\si{\volt}] \times CS_m [\si{\milli\ampere\hour\per\gram}]
\end{align*}
On peut alors calculer l'énergie théorique maximale, siglée ETM, pour chacune des électrodes :
\begin{center}
\begin{tabular}{ccccccc}
\toprule
Matériau & \ce{Li(s)} & \ce{LiC6(s)} & \ce{Li15Si4(s)} & \ce{LiCoO2(s)} & \ce{LiFePO4(s)} & \ce{LiNi_{1/3}Co_{1/3}Mn_{1/3}O2(s)} \\
\midrule
$f.e.m.$ [$\si{\volt}$] & 0 & 0.3 & 0.6 & 4.1 & 3.3 & 3.9  \\
Capacité [\si{\milli\ampere\hour\per\gram}] & 3860 & 372 & 3580 & 274 & 170 & 279 \\
ETM [\si{\milli\watt\hour\per\gram}] & $0$ & $112$ & $2148$ & $1123$ & $561$ & $1088$ \\
\bottomrule
\end{tabular}
\end{center}
La valeur d'un matériau comme électrode au pôle $\oplus$ doit tenir compte de deux choses : de la densité d'énergie stockée, mais aussi de la qualité de cette énergie.
}

\cteacher{Attention, ce diagramme est analogue au diagramme de Ragone mais n'en est pas un. Dans le diagramme de Ragone, on représente l'énergie en fonction de la puissance. Vous pouvez leur préciser cette erreur dans l'énoncé, bien qu'elle ne soit pas critique.}
Le diagramme de Ragone permet de comparer graphiquement les performances de dispositifs de stockage d'énergie (condensateurs, accumulateurs, pile à combustible...) ou de matériaux au sein d'un même dispositif, comme c'est le cas ici pour l'accumulateur au lithium.

\begin{figure}[H]
\subfloat[\label{fig:01a}]{\includegraphics[height=5cm]{Ex37/Fig01.png}}\subfloat[\label{fig:01b}]{\includegraphics[height=5cm]{Ex37/Fig02.png}}
\caption{Potentiel moyen de fonctionnement de quelques matériaux d'électrode négative et positive en fonction des valeurs de capacité accessibles expérimentalement ($<$ aux valeurs théoriques). Du coté anodique $\ominus$, les matériaux d’intercalation (spinelle LTO : \ce{Li4Ti5O12}, graphite, carbone dur), les matériaux de conversion (oxydes comme \ce{SnO2} et fluorures, phosphures et sulfures), les alliages (étain, germanium, silicium) et le lithium métallique. Du coté cathodique $\oplus$, les matériaux lamellaires [LTS : \ce{LiTiS2}, LCO : \ce{LiCoO2} et LCO partiellement substitués : NMC (\ce{Li(Ni,Mn,Co)O2}) et NCA (\ce{Li(Ni,Co,Al)O2})], les spinelles (LMO : \ce{LiMn2O4}, et LMO partiellement substitué LNMO) et les matériaux polyanioniques (olivine - LFP : \ce{LiFePO4}, LCP : \ce{LiCoPO4} et tavorite LFSF : \ce{LiFeSO4F})}
\end{figure}%http://reperageterrestre.free.fr/coordonnees.html


\Q Où trouve-t-on les matériaux de haute et basse densité d'énergie dans les graphes précédents ?
\solution{
Les matériaux de haute (resp. basse) densité d'énergie se trouve en haut (resp. bas) à droite (resp. gauche) du graphique de Ragone.
}
\Q Une grandeur est manquante pour caractériser pleinement une cellule électrochimique. Laquelle ? Comment pourrait-on l'y ajouter ?
\solution{
Bien que les batteries ne soient pas des dispositifs de puissance, les batteries doivent être capable de débiter une certaine énergie par unité de temps. \\
Pour représenter ces performances en puissance, on utilise le diagramme de Ragone qui représente l'énergie en fonction de la puissance, le plus souvent en échelle log-log :
\begin{align*}
\log ETM = \log P_s + \log \tau
\end{align*}
avec $\tau$ le temps caractéristique de décharge de l'accumulateur.
\begin{center}
\begin{minipage}{0.48\textwidth}
\centering
\includegraphics[width=0.5\textwidth]{Ex37/Corr01}
\end{minipage}\hfill\begin{minipage}{0.48\textwidth}
\centering
\includegraphics[width=0.5\textwidth]{Ex37/Corr02}
\end{minipage}
\end{center}
}